\begin{thebibliography}{99}

    \addcontentsline{toc}{chapter}{СПИСОК ИСПОЛЬЗОВАННЫХ ИСТОЧНИКОВ}
    
    \bibitem{1}
    Боресков, А. В. Графика трехмерных сцен : учебное пособие / А. В. Боресков, Е. В. Шикин. — М. : Диалог-МИФИ, 2005. — 704 с.
    
    \bibitem{2}
    Роджерс, Д. Алгоритмические основы машинной графики / Д. Роджерс ; пер. с англ. — М. : Мир, 1989. — 512 с.
    
    \bibitem{3}
    Ландсберг, Г. С. Оптика : учебное пособие для вузов / Г. С. Ландсберг. — 6-е изд., стереот. — М. : Физматлит, 2003. — 848 с.
    
    \bibitem{4}
    Shirley, P. Ray Tracing in One Weekend [Электронный ресурс]. — 2020. — Режим доступа: https://raytracing.github.io (дата обращения: 10.10.2025).
    
    \bibitem{5}
    DIALux evo. The software for smart lighting design [Электронный ресурс]. — Режим доступа: https://www.dialux.com/en-GB/dialux (дата обращения: 10.12.2025).
    
    \bibitem{6}
    Blender Documentation. Cycles Renderer [Электронный ресурс]. — Режим доступа: https://docs.blender.org/manual/en/latest/render/cycles/ (дата обращения: 11.12.2025).
    
    \bibitem{7}
    Pharr, M. Physically Based Rendering: From Theory To Implementation / M. Pharr, W. Jakob, G. Humphreys. — 3rd ed. — Morgan Kaufmann, 2016. — 1266 p.
    
    \bibitem{8}
    Unreal Engine 5. Lumen Technical Guide [Электронный ресурс]. — Режим доступа: https://docs.unrealengine.com/5.0/en-US/lumen-global-illumination-and-reflections-in-unreal-engine/ (дата обращения: 12.12.2025).
    
    \bibitem{9}
    Mitsuba 3: A Retargetable Geometric Laboratory [Электронный ресурс]. — Режим доступа: https://www.mitsuba-renderer.org/ (дата обращения: 12.12.2025).
    
    \bibitem{10}
    LuxCoreRender v2.6 Overview [Электронный ресурс]. — Режим доступа: https://luxcorerender.org/ (дата обращения: 13.12.2025).
    
    \bibitem{11}
    ISO/IEC 14882:2017. Programming languages — C++ [Электронный ресурс]. — Режим доступа: https://www.iso.org/standard/68735.html (дата обращения: 12.10.2025).
    
    \bibitem{12}
    Python 3.12 documentation [Электронный ресурс]. — Режим доступа: https://docs.python.org/3.12/ (дата обращения: 12.12.2025).
    
    \bibitem{13}
    pybind11 — Seamless operability between C++11 and Python [Электронный ресурс]. — Режим доступа: https://pybind11.readthedocs.io/ (дата обращения: 14.10.2025).
    
    \bibitem{14}
    NumPy Documentation [Электронный ресурс]. — Режим доступа: https://numpy.org/doc/ (дата обращения: 14.10.2025).
    
    \bibitem{15}
    Pillow (PIL Fork) Documentation [Электронный ресурс]. — Режим доступа: https://pillow.readthedocs.io/ (дата обращения: 15.10.2025).
    
    \bibitem{16}
    Pharr, M. Physically Based Rendering: From Theory To Implementation / M. Pharr, W. Jakob, G. Humphreys. — 4th ed. — MIT Press, 2023. — 1280 p.
    
    \bibitem{17}
    Marrs, A. Ray Tracing Gems II: Next Generation Real-Time Rendering with DXR, Vulkan, and OptiX / A. Marrs, P. Shirley, I. Wald. — Apress, 2021. — 720 p.
    
    \bibitem{18}
    Назаров, А. В. Компьютерная графика. Практикум : учебное пособие для вузов / А. В. Назаров, О. В. Назарова. — Санкт-Петербург : Лань, 2024. — 72 с.
    
    \bibitem{19}
    Ивлев, А. Н. Инженерная компьютерная графика : учебник / А. Н. Ивлев, О. В. Терновская. — 4-е изд., стер. — Санкт-Петербург : Лань, 2024. — 260 с.
    
    \bibitem{20}
    Грегори, М. Профессиональный C++ / М. Грегори ; пер. с англ. — 5-е изд. — М. : Диалектика, 2022. — 1328 с.
    
\end{thebibliography}